\section{Physical volume and marked balls}
\label{sec:volume-witnesses}

Write \(\mu_i=dV_{g_i}\), a probability measure, and fix
\(0<Y_0\leq Y_i\leq S\). The estimates in
Section~\ref{vol:physical-tests} require only the unit-volume Yamabe
minimizers and the scalar lower bound. The positive Ricci representative
in Section~\ref{vol:positive-ricci} uses the chiral upper bound as well.

\subsection{Cutoff tests on the physical source}
\label{vol:physical-tests}

For a Lipschitz cutoff \(0\leq\chi\leq1\), define
\[
 a=\int_M\chi^4\,d\mu_i,\qquad
 D=\int_M|d\chi|_{g_i}^2\,d\mu_i,\qquad
 A(\chi)=\left(\int_M|d\chi|_{g_i}^4\,d\mu_i\right)^{1/2}.
\]
If \(T\) contains \(\{0<\chi<1\}\) and the gradient support up to a
null set, the physical Yamabe inequality gives
\begin{align}
 \sqrt a-a&\leq 6D/Y_i+\mu_i(T)/4,\label{vol:yamabe-partition}\\
 Y_0(1-\sqrt{s})\sqrt m&\leq 6D
 \quad\text{if }\mu_i(\{\chi=1\})\geq m,\quad
                   \mu_i(\{\chi\ne0\})\leq s,\label{vol:small-support}\\
 D&\leq A(\chi)\sqrt{\mu_i(T)}.\label{vol:holder}
\end{align}
For the first inequality use
\(\int\chi^2\,d\mu_i\leq a+\mu_i(T)/4\).
For the second use \(\int\chi^2\,d\mu_i\leq\sqrt{s}\sqrt a\).
The last inequality is H\"older's inequality. The integral defining
\(A(\chi)\) is invariant under conformal changes of metric in dimension
four. A fixed cutoff in a uniformly bilipschitz comparison collar thus
has uniformly bounded \(A(\chi)\). If the comparison collar is a
length-\(L\) cylinder uniformly bilipschitz to the unit spherical product, a
cutoff with slope \(O(L^{-1})\) has \(A(\chi)=O(L^{-3/2})\).

The same Yamabe inequality yields a spectral test in the physical
metric. For smooth \(f\) with \(\int f\,d\mu_i=0\), expanding the
Yamabe quotient at \(1+t f\) to second order gives
\[
 6\int|df|_{g_i}^2\,d\mu_i\geq 2Y_i\int f^2\,d\mu_i .
\]
In particular \(\lambda_1(-\Delta_{g_i})\geq Y_i/3\), with
\(\Delta=\operatorname{div}\nabla\) and \(\lambda_1\) the first
positive eigenvalue. If \(\chi\) equals one and zero on disjoint sets
of masses \(m_1,m_0\), its variance is at least
\(m_1m_0/(m_1+m_0)\) for \(m_1+m_0>0\). Indeed, minimize
\(m_1(1-c)^2+m_0c^2\) over constants \(c\). Thus the spectral test and
\eqref{vol:holder} also bound any two such masses in terms of the
capacity and physical volume of their transition.

\begin{proposition}[A finite-exit volume witness]
\label{vol:exit-witness}
Let \(K_i\subset M\) be a marked compact domain with finitely many
complementary components \(U_{ij}\), \(1\leq j\leq J\), and null
boundaries. Suppose each \(U_{ij}\) has a cutoff equal to one there,
zero on the other components, with its transition in a compared
buffer \(T_{ij}\subset K_i\). Suppose \(A(\chi_{ij})\leq A\) and the
buffers have overlap multiplicity at most \(\ell\).
Write \(\kappa_i=\mu_i(K_i)\) and \(m_{ij}=\mu_i(U_{ij})\).
If \(m_{ij}\leq1-\eta\) for every \(j\), with \(0<\eta<1\), then
\begin{equation}
 \kappa_i\geq
 \min\left\{\frac{\eta}{2},\frac{c_\eta}{36A^2\ell+c_\eta}\right\},
 \qquad
 c_\eta=Y_0^2(1-\sqrt{1-\eta/2})^2 .
 \label{vol:exit-lower}
\end{equation}
If two exits each have mass at least \(\eta\), each of their buffers
has mass at least
\(Y_0^2(1-\sqrt{1-\eta})^2\eta/(36A^2)\) when \(A>0\).
If instead \(\kappa_i\to0\) at this fixed comparison stage, every limit
of the finite mass vector has exactly one entry equal to one.
\end{proposition}

\begin{proof}
For the first assertion, if \(\kappa_i<\eta/2\), the nonzero set of
each cutoff has mass at most \(m_{ij}+\kappa_i\leq1-\eta/2\).
Squaring \eqref{vol:small-support}, using \eqref{vol:holder}, and
summing gives
\(c_\eta(1-\kappa_i)\leq36A^2\ell\kappa_i\).
The case \(J=0\) has \(\kappa_i=1\).
For the second assertion, either cutoff has a nonzero set of mass
at most \(1-\eta\) and equals one on mass at least \(\eta\).
For the last assertion, at a limit of the mass vector the transition
terms in \eqref{vol:yamabe-partition} vanish by \eqref{vol:holder},
and \(a\to m_j\) for each cutoff. Hence \(\sqrt{m_j}\leq m_j\).
The masses sum to one. For successive coherent cuts with the same
exit labels, differences in exit masses are measured by the
intervening compared compact regions. If their masses tend to zero,
the surviving label agrees on all stages of a nested exhaustion
after diagonal subsequence extraction.
\end{proof}

For example, a bound \(\sup_x\mu_i(B_{g_i}(x,r))\leq1-\eta\) at a
fixed radius supplies \eqref{vol:exit-lower} with margin \(\eta/2\)
if each exit is contained in such a ball up to \(o(1)\) physical
volume. This application requires containment for the actual
complementary pieces.

\begin{proposition}[Saturation on a capacity-controlled exhaustion]
\label{vol:exhaustion}
Let \(X\) be an open comparison manifold with coherent embeddings
\(F_i\) on neighborhoods of exhausting compact domains. Let
\(\chi_b\in C_c^\infty(X;[0,1])\) equal one on those domains, with
compact transition supports \(T_b\) leaving every compact as
\(b\downarrow0\). Suppose the transported cutoffs, extended by zero to
\(M\), satisfy \(\sup_b\limsup_i A(\chi_b\circ F_i^{-1})<\infty\).
Then every locally vague limit \(\lambda\) of the physical measures
pulled back to \(X\) has total mass zero or one. A positive lower
mass bound on a fixed compared compact forces mass one. A positive
limsup does so along a subsequence realizing that limsup.
\end{proposition}

\begin{proof}
Vague convergence gives \(0\leq\lambda(X)\leq1\). For each fixed
\(b\), compactness of the transition support gives
\[
 \limsup_{i\to\infty}\mu_i(F_i(T_b))\leq\lambda(T_b).
\]
By \eqref{vol:holder}, the Dirichlet term vanishes as first
\(i\to\infty\), then \(b\downarrow0\). Taking those limits in the
Yamabe test for the transported cutoff gives
\(\sqrt{\lambda(X)}\leq\lambda(X)\): the compactly supported second
and fourth powers converge against the local measures for each \(b\),
and both integrals tend to \(\lambda(X)\) by dominated convergence.
The mass witness follows by testing with a compactly supported
continuous function equal to one on the witnessing compact.
For ordinary punctures of a smooth closed comparison, radial cutoffs
on annuli of radii \(b,2b\) give such a uniform fourth-gradient
bound after the source limit, by uniform local metric comparison
and conformal invariance.
\end{proof}

\begin{remark}[Evacuation in shrinking normal charts]
\label{vol:normal-chart-evacuation}
For any supplied smooth source sequence, comparison embeddings can
be chosen whose entire images have physical volume tending to zero.
Choose oriented source normal coordinates based at \(q'_i\in M\)
and radii \(R_i\to\infty\) in the parameter space. For each \(i\),
choose \(\lambda_i>0\) small enough that
\(F_i(x)=\exp_{q'_i}^{g_i}(\lambda_i x)\) is an embedding on
\(B_{R_i}(0)\), its image has \(\mu_i\)-mass at most \(i^{-1}\), and
\(\lambda_i^{-2}F_i^*g_i\) tends to the Euclidean metric smoothly on
compacts. Each such choice is possible by smoothness of the individual
metric in its normal coordinates. Under stereographic identification
with a growing subset of \(S^4\setminus\{p\}\), the comparison
\(m_i=\omega^2\lambda_i^{-2}F_i^*g_i\), where
\(\omega=2/(1+|x|^2)\), converges smoothly locally to the round metric.
Its physical factor is \(\lambda_i/\omega\), and every fixed compared
compact evacuates. This construction preserves all estimates on the
supplied source sequence. A source-only ball bound can certify
volume detection only with an additional relation to the marking.
\end{remark}

\subsection{Positive Ricci geometry at the half-volume scale}
\label{vol:positive-ricci}

The Schouten tensor here is
\(A_h=\tfrac12(\operatorname{Ric}_h-\tfrac16R_h h)\).
Signature and Gauss--Bonnet on the homotopy sphere give, in the
Hilbert--Schmidt curvature-operator normalization of \(E_+\),
\begin{equation}
 \int_M\sigma_2(A_h)\,dV_h=\frac{8\pi^2-E_+(g_i)}{2}
 \quad\text{for }h\in[g_i].
 \label{vol:sigma2}
\end{equation}
Indeed, \(E_-=E_+\), the full tensor Weyl integral is \(8E_+\),
and the Gauss--Bonnet integrand in these conventions is
\(\tfrac14|W|_{\rm tens}^2+4\sigma_2(A_h)\).
Also \(\sigma_2(A_h)=R_h^2/96-|\operatorname{Ric}_0(h)|^2/8\).
The unit-volume Yamabe representative therefore obeys
\(Y_i^2\geq48\delta\) when \(E_+(g_i)\leq8\pi^2-\delta\), so one
may take \(Y_0=\sqrt{48\delta}\).
The positive Yamabe sign and \eqref{vol:sigma2} give a smooth
positive-Ricci conformal representative
\(H_i^0\) by Chang--Gursky--Yang
\cite[Theorem A and Corollary B]{ChangGurskyYang2002PositiveRicci}.
Their Schouten tensor is \(2A_h\), which leaves the sign of its
second symmetric function unchanged. Any smooth positive-Ricci
representative may be used in the following normalization.

\begin{theorem}[Half-volume scale]
\label{vol:median}
Choose the least \(H_i^0\)-radius \(a_i\) of a ball carrying half of
\(\mu_i\), a witnessing center \(q_i\), and put
\(H_i=a_i^{-2}H_i^0\). In this subsection and the next, all unqualified
balls are taken in the \(H_i\) distance on the source or the
corresponding pointed distance on a limit. There are constants
\(C,\alpha>0\) depending only on \(Y_0\) such that
\begin{equation}
 \sup_x\mu_i(B_r(x))\leq C r^\alpha\quad(0<r\leq1),\qquad
 \mu_i(M\setminus B_R(q_i))\leq C R^{-\alpha}\quad(R\geq1).
 \label{vol:power}
\end{equation}
For every fixed \(D>0\) there is \(v_D>0\) such that
\begin{equation}
 v_D r^4\leq \operatorname{Vol}_{H_i}(B_r(x))
 \leq\omega_4 r^4\quad(x\in B_D(q_i),\ 0<r\leq1),\qquad
 \omega_4=\pi^2/2 .
 \label{vol:noncollapse}
\end{equation}
\end{theorem}

\begin{proof}
For each \(i\), distance spheres have zero Riemannian
four-volume by the polar-coordinate representation of measure.
The conformally related \(\mu_i\) also vanishes on every sphere.
Continuity of ball mass in center and radius follows from
inclusions between balls whose radii differ by the change in
center and radius. The maximum over the compact center set is
continuous in radius, starts at zero and reaches one. Thus
\(a_i>0\) is attained and, after scaling,
\(\mu_i(B_1(q_i))=\tfrac12=\max_x\mu_i(B_1(x))\).

For \(0<r<R\) use the logarithmic cutoff which equals one on
\(B_r(x)\), is affine in \(\log d_{H_i}(x,\cdot)\) on the annulus,
and vanishes outside \(B_R(x)\). Bishop's radial area bound in
nonnegative Ricci curvature supplies
\[
 \int_M |d\chi|_{H_i}^4\,dV_{H_i}
 \leq \frac{2\pi^2}{\log^4(R/r)}\int_r^R\frac{dt}{t}
 =\frac{2\pi^2}{\log^3(R/r)} .
\]
The reversed cutoff has the same cost. Put
\(C_0=72\pi^2/[Y_0^2(1-2^{-1/2})^2]\). When \(R\leq1\), apply
\eqref{vol:small-support} and \eqref{vol:holder} to the first
cutoff with \(s=1/2\) and transition mass
\(\mu_i(B_R(x))-\mu_i(B_r(x))\). When \(x=q_i\) and \(r\geq1\)
apply them to the reversed cutoff. This gives
\begin{align*}
 \mu_i(B_r(x))&\leq\frac{C_0}{C_0+\log^3(R/r)}\mu_i(B_R(x)),
       &&0<r<R\leq1,\\
 \mu_i(M\setminus B_R(q_i))&\leq
     \frac{C_0}{C_0+\log^3(R/r)}\mu_i(M\setminus B_r(q_i)),
       &&1\leq r<R .
\end{align*}
Fix \(L>0\), set \(\rho=C_0/(C_0+L^3)\),
\(\alpha=-\log\rho/L\), and iterate at ratios \(e^L\).
Monotonicity gives \eqref{vol:power} with \(C=(2\rho)^{-1}\).

By relative Bishop--Gromov comparison, a maximal \(r\)-separated
family with centers in \(B_D(q_i)\) has cardinality at most
\((1+4D/r)^4\). Its disjoint \(r/2\)-balls lie in
\(B_{D+r/2}(q_i)\). For any such center, the ball of radius
\(2D+r/2\) about that center contains this common enclosing ball.
Relative volume comparison at that center gives the claimed
packing bound. In particular a dimensionally bounded number of
\(1/4\)-balls covers \(B_1(q_i)\), so one of them has physical
mass at least \(c>0\), with \(c\) universal. Its concentric
\(1/2\)-ball has physical mass at most one half. The cutoff between
these balls has gradient bounded by \(4\) in \(H_i\).
Its fourth-gradient integral is bounded by a universal constant
times \(\operatorname{Vol}_{H_i}B_2(q_i)\).
Equations \eqref{vol:small-support} and \eqref{vol:holder} now give
\(\operatorname{Vol}_{H_i}B_2(q_i)\geq v>0\) independently of \(i\).
For \(x\in B_D(q_i)\), compare \(B_r(x)\) to \(B_{D+2}(x)\),
which contains the base ball. This yields the lower bound in
\eqref{vol:noncollapse}. The upper bound is Bishop comparison.
\end{proof}

\begin{corollary}[A protected ball on each source]
\label{vol:protected}
Fix \(A_0>1\), \(\zeta>0\), a nonnegative Borel measure \(\nu_i\) on
each source with \(\nu_i(M)\leq\Lambda<\infty\), and at most \(J\)
prescribed source points for each \(i\).
There exist \(D,r,m>0\), independent of \(i\), and centers
\(x_i\in B_D(q_i)\) such that
\[
 \mu_i(B_r(x_i))\geq m,\qquad
 \nu_i(B_{A_0r}(x_i))\leq\zeta .
\]
The centers may be chosen at distance at least \(2A_0r\) from every
prescribed point.
\end{corollary}

\begin{proof}
Take \(D\geq1\) so the tail bound leaves at most a quarter of the
mass outside \(B_D(q_i)\). For \(2A_0r\leq1\), Fubini gives
\[
 \int_M \nu_i(B_{A_0r}(x))\,d\mu_i(x)
 =\int_M\mu_i(B_{A_0r}(z))\,d\nu_i(z)
 \leq\Lambda C(A_0r)^\alpha .
\]
Choose \(r\) so \(C(2A_0r)^\alpha(J+\Lambda/\zeta)\leq1/4\).
By Markov's inequality and \eqref{vol:power}, centers in the desired
window with the desired buffer and mark clearance have physical
mass at least one half. A maximal \(r\)-separated family of these
centers has at most \((1+4D/r)^4\) members by the preceding
packing proof. Their \(r\)-balls cover the allowed centers.
Take \(m=\tfrac12(1+4D/r)^{-4}\).
\end{proof}

One can take \(\nu_i=|\operatorname{Rm}(g_i)|_{\rm tens}^2\,d\mu_i\).
In fact, the decomposition of the full curvature tensor and
\eqref{vol:sigma2} give
\[
 \int_M|\operatorname{Rm}(g_i)|_{\rm tens}^2\,d\mu_i
 =16E_+(g_i)-64\pi^2+Y_i^2/3\leq16E_+(g_i)+64\pi^2 .
\]
One can also take a chiral Weyl measure, using its conformal
invariance. These estimates control the indicated measures on
subsets of the source defined by \(H_i\)-distance.

\begin{proposition}[Bounded factor on bounded windows]
\label{vol:bounded-factor}
Write \(g_i=w_i^2H_i\). For each fixed \(D>0\),
\begin{equation}
 \sup_{B_D(q_i)} w_i\leq C_D,qquad
 d\mu_i\leq C_D^4\,dV_{H_i}\quad\text{on }B_D(q_i).
 \label{vol:factor-bound}
\end{equation}
In particular \(d_{g_i}(x,y)\leq C_{3D+1}d_{H_i}(x,y)\) for
\(x,y\in B_D(q_i)\).
\end{proposition}

\begin{proof}
All norms and gradients in this proof use \(H_i\).
Positive scalar curvature in this metric and the conformal
Yamabe equation give
\(-6\Delta_{H_i}w_i\leq Y_i w_i^3\).
For centers in \(B_{D+2}(q_i)\) and radii \(r\leq1\), the local
Sobolev inequality on a complete four-manifold with nonnegative
Ricci curvature \cite[Theorem 3.1]{SaloffCoste1992Uniformly}
and \eqref{vol:noncollapse} give
\[
 \|z\|_4^2\leq C_D\bigl(\|dz\|_2^2+r^{-2}\|z\|_2^2\bigr),
 \qquad z\in H_0^1(B_r(x)).
\]
The relative prefactor in that theorem is bounded by
\(C r^2\operatorname{Vol}_{H_i}(B_r(x))^{-1/2}\).
For \(p\geq2\) test the subsolution with \(\eta^2w_i^{p-1}\)
for a cutoff supported in the ball. The principal term is
\(24(p-1)p^{-2}\int\eta^2|d(w_i^{p/2})|^2\). Absorbing half of
it against the cutoff cross term and using the Sobolev inequality
for \(z=\eta w_i^{p/2}\) yields
\begin{equation}
 \|z\|_4^2\leq C_Dp^2\int (|d\eta|^2+r^{-2}\eta^2)w_i^p\,dV_{H_i}
       +C_Dp^2\int Y_i w_i^2z^2\,dV_{H_i}.
 \label{vol:power-test}
\end{equation}
At \(p=4\), the last integral is at most
\(S\mu_i(B_r(x))^{1/2}\|z\|_4^2\). By \eqref{vol:power}, take a
fixed \(r=r(D)>0\) for which the term absorbs. A cutoff equal to
one on \(B_{r/2}(x)\) and \(\int w_i^4\,dV_{H_i}=1\) give a
uniform \(L^8\) bound there.

Thus \(V_i=Y_iw_i^2\) is uniformly bounded in \(L^4\) on the
smaller ball. For subsequent cutoffs,
\[
 \int V_i z^2\,dV_{H_i}\leq C_{D,r}\|z\|_{8/3}^2
       \leq C_{D,r}\|z\|_2\|z\|_4 .
\]
Young's inequality in \eqref{vol:power-test} absorbs the \(L^4\)
factor, with coefficient at most \(C_{D,r}p^4\) on the squared
\(L^2\) norm of \(z\). Iterate on balls decreasing from radius
\(r/2\) to \(r/4\), with derivative bounds \(C2^j/r\) and
\(p_j=4\cdot 2^j\). The successive norm multipliers are bounded by
\([C_{D,r}p_j^4 4^j]^{1/p_j}\), whose product is finite. The
centers range over the bounded window. This proves the first
estimate in \eqref{vol:factor-bound}, and the second is the
four-dimensional change of volume. An \(H_i\)-minimizing geodesic
between the indicated points remains in \(B_{3D+1}(q_i)\), so
integrating \(w_i\) along it gives the distance estimate.
\end{proof}

\subsection{Regular points and smooth marked balls}

\begin{theorem}[Physical mass at metric regular points]
\label{vol:regular}
After passing to a subsequence, \((M,d_{H_i},q_i)\) converges
in pointed Gromov--Hausdorff distance to a proper metric space
\((X,d,q)\). In a joint pointed realization, \(dV_{H_i}\) tends
locally weakly to Euclidean-normalized \(\mathcal H^4\) on \(X\),
and \(\mu_i\) converges to a probability
\(\mu=f\,\mathcal H^4\) with \(f\in L^\infty_{\rm loc}(X)\).
The probability \(\mu\) gives full mass to points whose tangent
cones are all Euclidean four-space. If \(\nu_i\) are as in
Corollary~\ref{vol:protected}, take a local weak limit \(\nu\)
on the same subsequence. For any fixed \(A_0>1\), one can choose a
regular point \(x\in\operatorname{supp}\mu\), off any prescribed
finite set and the atoms of \(\nu\), with approximating source
centers \(x_i\to x\) and
\begin{equation}
 \lim_{b\downarrow0}\lim_i\frac{\mu_i(B_b(x_i))}{\omega_4 b^4}
      = f(x)\in(0,\infty),\qquad
 \lim_{b\downarrow0}\lim_i\nu_i(B_{A_0b}(x_i))=0 .
 \label{vol:regular-iterated}
\end{equation}
The outer limit is through joint continuity radii for the
respective measures.
\end{theorem}

\begin{proof}
The packing bound gives pointed metric precompactness.
Noncollapse at \(q_i\) and \(\operatorname{Ric}_{H_i}\geq0\)
allow application of Cheeger--Colding's noncollapsed volume
convergence theorem
\cite[Theorem 5.9]{CheegerColding1997StructureI}.
Local mass bounds give joint local weak limits for the physical
probabilities and for \(\nu_i\) after further subsequence
extraction. One can use a common pointed metric realization
with convergent centers on each bounded window. The physical
tail bound in \eqref{vol:power} preserves all the probability.
Testing \eqref{vol:factor-bound} against nonnegative continuous
functions of compact support yields the density assertion.
The same volume convergence and the regularity theorem
\cite[Theorems 2.1 and 5.11]{CheegerColding1997StructureI}
give Euclidean four-dimensional tangent cones at
\(\mathcal H^4\)-almost every point, hence at \(\mu\)-almost
every point. At such a regular point,
\(\mathcal H^4(B_b(x))/(\omega_4 b^4)\to1\).
This follows from noncollapsed volume convergence applied to
the Euclidean blow-ups, using a diagonal sequence of the
smooth sources at those scales.

The measure \(\mathcal H^4\) is locally doubling by the limiting
Bishop--Gromov comparison. Differentiation of \(\mu\) relative
to \(\mathcal H^4\) gives the ball density ratio \(f(x)\) at
\(\mu\)-almost every point, where \(0<f(x)<\infty\).
The at most countably many atoms of \(\nu\) and any prescribed
finite set are \(\mu\)-null. Choose a point with all these
properties. First transfer the two balls at fixed radii using
local weak convergence and their continuity boundaries.
Then take \(b\downarrow0\), using differentiation and
\(\nu(\{x\})=0\).
\end{proof}

\begin{corollary}[Finitely many disjoint smooth source balls]
\label{vol:smooth-balls}
Fix \(0<\epsilon<1\), \(\zeta>0\), and a maximum radius
\(r_*>0\). On the subsequence above, for all sufficiently large
\(i\) there are uniformly finitely many disjoint smoothly
embedded closed standard four-balls \(\Omega_{i\ell}\subset M\).
Each lies inside an \(H_i\)-ball of radius at most \(r_*\),
has \(\nu_i(\Omega_{i\ell})<\zeta\), and the union has physical
mass at least \(1-\epsilon\). The embeddings can be oriented
consistently with \(M\).
\end{corollary}

\begin{proof}
Choose a bounded window with physical limiting mass greater than
\(1-\epsilon/4\). Fix \(\tau\in(0,1/6)\) sufficiently small that
\((1-3\tau)^4(1-\epsilon/2)>1-\epsilon\), with strict room to
decrease the left side slightly. At almost every point for \(\mu\)
in the window, the preceding proof provides both Euclidean
\(\mathcal H^4\)-density and positive physical four-density, with
zero atomic \(\nu\)-mass. Fix a volume error small enough that
the almost-maximal-volume ball theorem
\cite[Appendix A, Theorems A.1.8--A.1.9]{CheegerColding1997StructureI}
provides a smooth open-ball embedding into any qualifying source
ball, with image containing the concentric ball of relative
radius \(1-\tau\). That theorem applies to the smooth sources
with \(\operatorname{Ric}\geq0\) at sufficiently small radii.

Use a fine family of limit balls about the indicated points,
with radii \(r<r_*\) below this theorem's bound, and
\(\nu(\overline B_r)<\zeta/2\). The radii
\(r,(1-\tau)r,(1-3\tau)r\) may all be taken as continuity radii.
Choose them sufficiently small that the \(\mathcal H^4\)-volume
ratio at \((1-\tau)r\) meets the theorem's volume requirement
with a strict margin. At the same time the physical mass ratio
between \((1-3\tau)r\) and \(r\) can be made as close to
\((1-3\tau)^4\) as desired. The Vitali covering theorem for the
locally doubling \(\mathcal H^4\) measure selects disjoint balls
covering almost all the indicated centers. Since
\(\mu\ll\mathcal H^4\), a finite subfamily covers physical
mass greater than \(1-\epsilon/2\). By the ratio choice its
\((1-3\tau)r\)-balls still carry strictly more than
\(1-\epsilon\). Shrinking the selected outer balls to
\((1-\tau)r\) gives a positive separation margin for this
finite family.

With all radii held fixed, transfer the centers to the sources
using the pointed realization. Volume convergence and the
strict margin allow application of the smooth ball theorem
inside each \((1-\tau)r\)-ball for large \(i\). Its image
contains the closed concentric \((1-3\tau)r\)-ball compactly,
since \((1-\tau)^2>1-3\tau\). The preimage of this smaller
metric ball is compact in the open parameter ball, so it is
contained in the interior of a smaller closed concentric
parameter ball. Take its image for \(\Omega_{i\ell}\).
Disjointness follows from the separation margin. Convergence
of the measures on the fixed smaller and outer balls gives
the total physical mass and the per-ball bound on \(\nu_i\).
A reflection of a parameter ball, if necessary, fixes its
orientation. Each boundary is an actual smooth three-sphere.
\end{proof}

The ball theorem here supplies smooth embeddings at each index.
The estimates concern their volumes and the specified measures.
First derivatives of the marked metrics in the critical norm and
a comparison on a whole punctured truncation require further
geometric input.

\subsection{Saturation for a whole marking at the same scale}

\begin{proposition}[Exits in the half-volume gauge]
\label{vol:normalized-marking}
Suppose there are coherent whole-truncation markings
\(\Phi_i\) of a connected smooth closed \((N,h)\) minus a
finite set \(P\) of ordinary points, such that
\(\Phi_i^*H_i\to h\) locally uniformly. Write
\(K_b=N\setminus\bigcup_{a\in P}B_h(a,b)\) for the compact
truncations by disjoint small normal balls. Then
\begin{equation}
 \lim_{b\downarrow0}\limsup_i
 \mu_i\bigl(M\setminus\Phi_i(K_b)\bigr)=0 .
 \label{vol:exit-saturation}
\end{equation}
If the marking also has the local first-order and comparison
curvature properties of Definition~\ref{def:marking} with
\(m_i=\Phi_i^*H_i\), its conformal factors have strong local
\(H^1\) and \(L^4\) convergence and a positive smooth unit-mass
limit, including when \(y=S\).
\end{proposition}

\begin{proof}
Each spherical face of the connected marked core is a separating
hypersurface in \(M\), since \(H^1(M;\mathbb Z)=0\). The connected
core lies on its inward side. All other faces lie on that same
side, so each component of the complement meets exactly one face:
it cannot connect the outward side of one face to the outward
side of another while avoiding the core. If \(P=\varnothing\),
the embedding of closed \(N\) is onto by invariance of domain
and compactness, and there is no exterior.

For a mark \(a\), fix a small \(s>0\) less than half the separation
of the normal neighborhoods of distinct marks, and choose
\(b\ll s\). Take any point \(z_i\) of the face
\(\Sigma_i=\Phi_i(S_h(a,b))\). Its intrinsic \(H_i\)-diameter
is at most \(C_1 b\) for large \(i\). Following a face path and
then a normal radial segment in the core shows that the marked
annulus \(b<d_h(a,\cdot)<s-C_2 b\), with \(C_2>C_1+1\), lies
in \(B_s(z_i)\) for large \(i\) at fixed \(b,s\).
Its \(H_i\)-volume is at least
\[
 \omega_4s^4\bigl(1-O(s^2)-O(b/s)-o_i(1;b,s)\bigr).
\]
If a point of the exterior reaches depth \(s/2\) from \(\Sigma_i\),
choose a point \(x_i\) of exact depth \(s/2\) by continuity
along an exterior path to the face. Then
\(d_{H_i}(x_i,z_i)\leq s/2+C_1 b\). The ball \(B_{s/8}(x_i)\)
lies entirely in the exterior and in \(B_s(z_i)\).
The larger ball \(B_{2s}(x_i)\) contains the marked annulus.
Relative Bishop comparison at \(x_i\) implies that the exterior
ball has \(H_i\)-volume at least \(16^{-4}\) times that of
the annulus. The exterior ball and annulus are disjoint.
Choose \(s\) small, then \(b/s\) small, then \(i\) large.
Their combined volume would exceed \(\omega_4 s^4\), contradicting
Bishop comparison at \(z_i\). Therefore every exterior point
has depth less than \(s/2\), and the entire exterior is contained
in \(B_s(z_i)\).

By \eqref{vol:power}, for each sufficiently small \(s\), all
sufficiently small fixed \(b\) have limiting upper exterior mass at
most \(|P|Cs^\alpha\). This proves \eqref{vol:exit-saturation}
in the stated order of limits.

For the factor assertion, a sufficiently small fixed coordinate
ball on the marked core maps into a comparably small \(H_i\)-ball
by the uniform local metric comparison and lengths of coordinate
segments. Its \(L^4\) physical factor mass is therefore uniformly
small by \eqref{vol:power}. In the local critical-absorption estimate
of Proposition~\ref{prop:factor}, the difference of the factors
from their weak limit has small \(L^4\) norm there by fourth-power
splitting. The coefficient of its pure cubic term,
\(Y_i\|u_i-u\|_4^2\), is then below the local Sobolev coefficient
even for \(y=S\). The same absorption proof supplies strong local
\(H^1\) and \(L^4\) convergence. Equation~\eqref{vol:exit-saturation}
now gives unit mass for the limiting factor. The removable-point
and positivity argument of Proposition~\ref{prop:factor} then
gives the smooth positive factor on \(N\).
\end{proof}
